This study presents an aggregated analysis on the evolution of armed conflict in Mexico. The criminal war in Mexico is extremely complex: Drug Trafficking Organizations vie for control against one another and against government actors. Civilians are often subjected to DTO-related violence and left unprotected by the state. This cycle of violence is characterized by a set of interlocking interactions between armed actors. Existing research, however, typically ignores the interdependencies inherent to these networks.  Using a new collection of machine-coded event data, we generate conflict networks representing each year from 2004 to 2010.  Importantly, we demonstrate how these networks capture the independent nature of the Mexican conflict and test an important theoretical question of how protests influence the probability of violence between armed actors within the network.  Finally we employ a latent space approach to uncover previously unobservable violence between government actors and criminal groups. 
